\chapter{Your first engrave}
\label{ch:engrave}

Cutting is one-dimensional: the beam follows a line. Engraving removes material
from an \emph{area}, and that makes it a different craft with different dials ---
resolution, direction, dot patterns. This chapter covers both flavours of it:

\begin{itemize}
  \item \textbf{raster engraving}, where MeerK40t converts vector artwork into
        scan lines and burns it, and
  \item \textbf{image engraving}, where the artwork is a picture and the pixel
        values decide how hard the beam burns.
\end{itemize}

\noindent Both end up as scan lines at the machine. The difference is where the
decisions are made.

\begin{figure}[htbp]
\centering
\begin{tikzpicture}[scale=1.0]
\draw[fill=mklight,draw=mkgrey!60] (0,0) rectangle (7.2,2.2);
\foreach \y in {0.2,0.5,0.8,1.1,1.4,1.7,2.0} {
  \draw[mkred,line width=0.7pt] (1.0,\y) -- (6.2,\y);
  \draw[mkblue,line width=0.5pt,-{Stealth[length=1.4mm]}] (6.2,\y) -- (6.6,\y);
  \draw[mkblue,line width=0.5pt,-{Stealth[length=1.4mm]}] (1.0,\y) -- (0.6,\y);
}
\draw[{Stealth[length=2mm]}-{Stealth[length=2mm]},mkdark] (0.6,-0.35) -- (1.0,-0.35);
\node[font=\sffamily\scriptsize,color=mkdark] at (0.0,-0.35) {overscan};
\draw[{Stealth[length=2mm]}-{Stealth[length=2mm]},mkdark] (1.0,2.55) -- (6.2,2.55);
\node[font=\sffamily\scriptsize,color=mkdark,fill=white] at (3.6,2.55)
  {scan line, burned left to right and right to left};
\draw[{Stealth[length=2mm]}-{Stealth[length=2mm]},mkdark] (7.6,0.2) -- (7.6,0.5);
\node[font=\sffamily\scriptsize,color=mkdark,anchor=west] at (7.7,0.35) {step (DPI)};
\end{tikzpicture}
\caption{How a raster is burned. The head sweeps across the shape, reverses, and
steps down one line at a time. \emph{Overscan} is the extra run-up at each end
that lets the head reach full speed before the beam starts; without it, the edges
of the raster are darker than the middle. \emph{Directional Raster} decides
whether the return pass also burns. The line spacing is set by DPI.}
\label{fig:raster}
\end{figure}

\section{Raster: filling vector artwork}

Use this when the thing you want is a \emph{shape} --- a filled logo, a block of
text, a data plate --- and you want it engraved rather than cut. MeerK40t renders
the assigned elements to an image internally and scans it out.

\begin{walkthrough}{Engrave a filled logo or text block}
\begin{enumerate}[label=\arabic*.,leftmargin=1.4em]
  \item Draw or import the artwork. Filled shapes work best; a thin outline
        engraves as a thin outline.
  \item Right-click \emph{Operations} in the tree and choose \menu{Append
        operation>Append Raster}. A black operation appears.
  \item Select the artwork and assign it --- \key{Ctrl+A} then click the black
        square in the \panel{Operations} pane, or right-click the elements and
        use \menu{Assign Operation}.
  \item Open \menu{Tools>Properties} on the raster operation and set the two dials
        that matter first: \btn{DPI} and \btn{Speed}.
  \item Add \btn{Passes} if one pass is too pale; a second pass doubles the time
        and usually darkens more evenly than raising power does.
\end{enumerate}
\end{walkthrough}

\subsection*{Choosing DPI}

DPI and the machine's spot size are the same conversation. If your spot is about
0.1\unit{mm}, then your step is 0.1\unit{mm}, and there is no point scanning at
more than about 250\unit{dpi}, because the next line lands where the last one
burned. In practice:

\begin{center}\small
\begin{tabular}{@{}L{1.9cm}L{2.6cm}L{3.1cm}L{6.7cm}@{}}
\toprule
\textbf{DPI} & \textbf{Step} & \textbf{Typical use} & \textbf{Notes} \\
\midrule
100 & 0.254\unit{mm} & quick marking, deep text & visible lines; fast \\
250 & 0.102\unit{mm} & general CO\textsubscript{2} engraving &
    matches a normal focus spot; a good default \\
500 & 0.051\unit{mm} & fine detail, photographs & the shipped default; twice
    the time of 250 for often the same practical darkness \\
1000 & 0.025\unit{mm} & diode machines with tiny spots & slow; check that the
    result is actually finer before keeping it \\
\bottomrule
\end{tabular}
\end{center}

\noindent Time scales with the \emph{area} you fill and inversely with the line
step and the speed: doubling DPI doubles the number of lines, halving the area of
the same shape cuts the time in half. The \panel{Simulation} window's time
estimate will tell you the number exactly; the arithmetic above tells you which
way to push.

\subsection*{Direction, overscan and the rest}

\begin{itemize}
  \item \btn{Direction} (top-to-bottom, bottom-to-top, left/right, crosshatch,
        greedy, crossover, spiral, diagonal) changes the grain of the fill. For
        wood, a crosshatch looks woven; a single direction looks brushed.
  \item \btn{Directional Raster} on means the return pass also burns, roughly
        halving the time. Turn it off if your machine shows backlash at line ends
        --- the classic symptom is a faint ghosting to one side of each line.
  \item \btn{Overscan} (default 0.5\unit{mm}) is the run-up described in
        Figure~\ref{fig:raster}. If the first and last few millimetres of your
        engraving are noticeably darker, increase it. Ruida machines are reported
        to need it at zero, so check the device notes in Appendix~\ref{app:drivers}.
  \item \btn{Use grayscale instead} and \btn{Consider laserdot} are refinement
        switches: the first rasterises the artwork's antialiasing, the second
        skips pixels the beam already covered.
  \item \btn{Start Preference} decides which corner the scan begins in. On
        machines that are slow to reach speed, starting from the corner nearest
        the previous operation saves real seconds.
\end{itemize}

\section{Image: engraving a picture}

Photographs and logos supplied as \file{.png}, \file{.jpg} or another bitmap
format become image elements, and image elements are burned by an
\mkseries{Image} operation. The image's own properties carry the dithering and
tone controls, so the workflow is: import, adjust the picture, then engrave.

\subsection*{Importing}

\begin{walkthrough}{Bring in a photograph}
\begin{enumerate}[label=\arabic*.,leftmargin=1.4em]
  \item \menu{File>Import File}, choose your picture. It lands in the elements
        branch, selected.
  \item Check its size. Preferences (\tabname{Input/Output}) control what happens
        on import: \btn{Image DPI Scaling} decides whether the picture's stored
        DPI is trusted for physical size (unset: treat it as 96\unit{dpi});
        \btn{Center image on import} places it in the middle of the bed instead
        of the top-left corner; \btn{Scale oversized images} shrinks anything
        bigger than the bed; \btn{Create a file-node for imported image} wraps it
        in a file node for tidier trees.
  \item Resize deliberately. Set the real width in the \panel{Position} pane
        rather than dragging --- you usually know how wide the engraved picture
        should be, and DPI is derived from it.
  \item Right-click \emph{Operations}, \menu{Append operation>Append Image}.
        Assign the image to it.
  \item Open \menu{Tools>Properties} on the image element (or right-click it in
        the tree and choose \menu{Image properties}) to reach the panels that
        matter: \btn{DPI}, \btn{Dither}, \btn{Contrast}, \btn{Halftone},
        \btn{Tone}, \btn{Sharpen}, \btn{Gamma}, \btn{Edge Enhance},
        \btn{Auto Contrast}, and the \btn{Modification} script list.
\end{enumerate}
\end{walkthrough}

\subsection*{Dither: how a photograph becomes laser pulses}

A laser is either on or off. A photograph has thousands of shades. Dithering is
the trick that bridges them: shades become patterns of dots, dense where the
picture is dark and sparse where it is light. The pattern you choose changes the
character of the result as much as its darkness.

\begin{center}\small
\begin{tabular}{@{}L{4.3cm}L{9.9cm}@{}}
\toprule
\textbf{Dither algorithm} & \textbf{Character} \\
\midrule
\emph{Floyd-Steinberg} (default) & Error-diffusion classic. Smooth gradients,
fine grain. The safe default for photographs. \\
\emph{Legacy-Floyd-Steinberg} & The older implementation; use it to reproduce
work made before 0.9. \\
\emph{Atkinson} & Softer, more contrast, less error carried --- the old
Macintosh look. Good for portraits on wood. \\
\emph{Jarvis-Judice-Ninke}, \emph{Stucki}, \emph{Burkes}, \emph{Sierra3},
\emph{Sierra2}, \emph{Sierra-2-4a}, \emph{Shiau-Fan}, \emph{Shiau-Fan-2} &
Variants of error diffusion with progressively sharper or smoother character.
Try two or three on a test piece rather than over-thinking the names. \\
\emph{Bayer}, \emph{Bayer-Blue} & Ordered dithering: a regular grid of dots.
Looks intentional and graphic --- excellent for graphics, logos and gradients you
want to look mechanical, poor for photographs. \\
\bottomrule
\end{tabular}
\end{center}

\noindent Turn \btn{Dither} off if the image is already pure black and white (a
scanned line drawing, for instance) --- there is nothing to dither, and the result
is cleaner without it.

\subsection*{Tone controls}

The panels under the image properties do what the equivalent commands in an image
editor do, and they apply to the laser output rather than to a screen:

\begin{center}\small
\begin{tabular}{@{}L{3.0cm}L{11.2cm}@{}}
\toprule
\textbf{Panel} & \textbf{What it does} \\
\midrule
\btn{Contrast} & \emph{Contrast Amount} ($-127$\dots127) and \emph{Brightness
Amount}. The first thing to reach for on a flat-looking photo. \\
\btn{Halftone} & Converts to dot-screen style output with \emph{Sample},
\emph{Angle} and \emph{Oversample} sliders and a \emph{Use black rather than
white dots} switch. This is how you get a newspaper-print look. \\
\btn{Tone} & A full 256-point tone curve. Precise and slow to use; the right tool
when a specific band of the picture must burn darker. \\
\btn{Gamma} & Shifts midtones with a single factor. The most useful control for
matching a photo to a material's burn response. \\
\btn{Sharpen}, \btn{Edge Enhance} & Emphasise edges. Useful on wood, where
subtle detail is lost in the grain. \\
\btn{Auto Contrast} & A one-click stretch of the tonal range. \\
\btn{Modification} & Applies a named \emph{Raster-Wizard} script ---
\emph{Gold}, \emph{Stipo}, \emph{Gravy}, \emph{Xin}, \emph{Newsy} or
\emph{Simple} --- as an ordered list of image operations you can edit. \\
\btn{Crop} & Crop by pixel, with four sliders. \\
\bottomrule
\end{tabular}
\end{center}

\noindent If the picture needs more than this --- retouching, colour removal,
serious sharpening --- do it in an image editor first. MeerK40t's job is to burn
the picture well, not to be Photoshop.

\subsection*{Grayscale and depth maps}

For materials that respond to dwell rather than to dots --- and for diode lasers
where power modulation is smooth --- the image can be burned as \emph{grayscale}
instead of dithered: the pixel's brightness sets the power. The
\btn{Depthmap} checkbox goes further: it treats every grey level as a number of
times that pixel will be burned, with a selectable resolution from \emph{256 -
8bit} down to \emph{4 - 2bit}. Depth maps are ideal for engraving height into
wood, and they are slow, because repeated burns mean repeated passes over the
same area. The \btn{Grayscale} box also offers \emph{Invert} and
\emph{Red}/\emph{Green}/\emph{Blue}/\emph{Lightness} sliders --- the last is the
quickest way to lift a dark photograph before engraving.

\section{Big images, tiles and keyholes}

Two tools exist for the practical problems of large or awkward artwork, both
under \panel{Image ops} (\menu{Tools>Editing>Image Splitting}):

\begin{itemize}
  \item \textbf{Create split images} (\tabname{Render + Split}) renders the
        selection to an image at a chosen \textbf{DPI} and cuts it into an
        \emph{X-Axis} $\times$ \emph{Y-Axis} grid (up to 25 by 25). Each tile
        becomes its own element, which is what you want when a pattern is wider
        than the bed: engrave the tiles in separate runs, moving the material or
        using the machine's own repeat feature.
  \item \textbf{Keyhole operation} (\tabname{Keyhole operation}) renders the
        chosen object as a mask at a given \textbf{DPI}, with options to
        \emph{Invert} the mask and to \emph{Trace} the keyhole outline. This is
        how you engrave one shape \emph{through} another --- a picture inside a
        frame, or a logo that must avoid a pre-existing hole.
\end{itemize}

\noindent The same two operations are available in the console as
\cmd{render_split <cols> <rows> <dpi>} and \cmd{render_keyhole <dpi>}, which is
the faster route once you know your numbers.

\begin{machinebox}{Engraving is where fires start}
Raster and image jobs concentrate energy in one small area for minutes at a time,
on the same spot. Three habits prevent trouble: keep air assist on, keep the
extraction running, and do not walk away. If a job smells like a bonfire rather
than like wood, stop it and lower the power --- the smell of a raster job that is
too aggressive is unmistakeable once you have met it.
\end{machinebox}

\begin{tryit}{Three-pass test on the same picture}
The fastest way to learn the tone controls is to burn the same small photograph
three times, changing one thing each time.
\begin{enumerate}[label=\arabic*.,leftmargin=1.4em]
  \item Import a photograph roughly 40\unit{mm} wide; convert it to grayscale and
        raise its contrast in an image editor first.
  \item Duplicate it twice (\menu{Edit>Copy}, \menu{Edit>Paste}) and line the
        three copies up side by side with the \panel{Alignment} window
        (\btn{Align} tab, \emph{Relative to: Selection}, \emph{Treatment: As
        Group}).
  \item Give all three copies one \mkseries{Image} operation at 250\unit{dpi},
        then set different dither algorithms on each element ---
        \emph{Floyd-Steinberg}, \emph{Atkinson} and \emph{Bayer}.
  \item Engrave once and compare. Repeat with \emph{DPI} 250 versus 500 on the
        best-looking setting.
  \item Keep the sample in your material notebook and write the settings on the
        back. The next photo job takes one minute instead of thirty.
\end{enumerate}
\end{tryit}
